\section{绪论}
\subsection{研究背景与意义}
随着深度学习技术的快速发展，神经网络在图像分类、目标检测等任务中取得了优异效果。但现实场景中数据集普遍存在类别不均衡问题，模型训练时容易偏向样本数量较多的大类，在少数类别上识别精度偏低，泛化能力下降。

波动力学从波动、演化的角度描述状态传播规律，和神经网络特征逐层传递、梯度迭代更新的过程存在相似性。将类别偏差类比为波动信号，可以用来刻画偏置随训练不断放大的过程，为预测、抑制类别偏置提供新的思路。

本研究旨在构建结合波动力学思想的类别偏置预测模型，在训练过程中提前预测各类别偏差演化趋势，动态调整损失权重，缓解不均衡分类下模型偏向大类的问题，提升模型整体均衡分类性能。

\subsection{国内外研究现状}
在类别不均衡与类别偏置方向，现有方案主要分为数据层面与算法层面：数据层面包含重采样、生成样本等方式；算法层面包含加权损失、两阶段训练等策略。这类方法大多属于事后补偿手段，缺少对偏置演化趋势的预测机制。

波动力学相关思想在机器学习领域多用于图信号传播、图神经网络等方向，将波动方程用于模拟信号在节点之间的传播、衰减过程。目前较少有研究将波动演化思想用于建模分类任务中的类别偏置动态变化，因此该方向具备研究价值。

\subsection{研究内容与技术路线}
本文主要研究内容如下：
\begin{enumerate}
    \item 分析不均衡数据集下类别偏置产生机理，选取合适指标量化各类别偏置大小；
    \item 借鉴波动力学波动传播思想，构建波动偏置预测模块，建模偏差随训练迭代的传播过程；
    \item 设计分类损失与偏置预测损失联合优化方案，完成模型训练；
    \item 在公开不均衡数据集上开展对比实验，验证模型对类别偏置的抑制效果。
\end{enumerate}

技术路线：数据集构建与偏置指标定义 → 波动特征建模与预测模块设计 → 模型联合损失构建与训练 → 验证集调参、测试集对比分析。

\subsection{论文组织结构}
第一章为绪论，介绍研究背景、意义、国内外现状、研究内容与论文框架；
第二章介绍相关理论基础，包括类别不均衡、神经网络特征提取、波动基础理论以及不均衡学习经典方法；
第三章提出基于波动力学的类别偏置预测模型；
第四章设计实验方案，对比基线模型并分析实验结果；
第五章总结全文，指出不足与后续改进方向。
